\documentclass[11pt]{article}
\usepackage[utf8]{inputenc}
\usepackage[T1,T2A]{fontenc}
\usepackage[russian]{babel}
% символы Unicode исходного текста Lean, чтобы цитировать его дословно
\DeclareUnicodeCharacter{2286}{\ensuremath{\subseteq}}
\DeclareUnicodeCharacter{2192}{\ensuremath{\to}}
\DeclareUnicodeCharacter{211D}{\ensuremath{\mathbb{R}}}
\DeclareUnicodeCharacter{2200}{\ensuremath{\forall}}
\DeclareUnicodeCharacter{2203}{\ensuremath{\exists}}
\DeclareUnicodeCharacter{2264}{\ensuremath{\le}}
\DeclareUnicodeCharacter{2227}{\ensuremath{\wedge}}
\DeclareUnicodeCharacter{2243}{\ensuremath{\simeq}}
\DeclareUnicodeCharacter{1D62}{\ensuremath{_i}}
\DeclareUnicodeCharacter{2194}{\ensuremath{\leftrightarrow}}
\DeclareUnicodeCharacter{2260}{\ensuremath{\ne}}
\DeclareUnicodeCharacter{2211}{\ensuremath{\sum}}
\DeclareUnicodeCharacter{2208}{\ensuremath{\in}}
\DeclareUnicodeCharacter{222A}{\ensuremath{\cup}}
\DeclareUnicodeCharacter{221A}{\ensuremath{\surd}}
\DeclareUnicodeCharacter{2216}{\ensuremath{\setminus}}
\DeclareUnicodeCharacter{2115}{\ensuremath{\mathbb{N}}}
\usepackage{amsmath,amssymb,amsthm,mathtools}
\usepackage{tikz}
\usepackage{upquote}
\usepackage{booktabs}
\usepackage[margin=1in]{geometry}
\setlength{\emergencystretch}{3em}
\usepackage[colorlinks=true,linkcolor=blue,citecolor=blue,urlcolor=blue]{hyperref}

\newtheorem{theorem}{Теорема}[section]
\newtheorem{lemma}[theorem]{Лемма}
\newtheorem{proposition}[theorem]{Предложение}
\newtheorem{corollary}[theorem]{Следствие}
\theoremstyle{definition}
\newtheorem{definition}[theorem]{Определение}
\theoremstyle{remark}
\newtheorem{remark}[theorem]{Замечание}
\newtheorem{example}[theorem]{Пример}

\newcommand{\R}{\mathbb{R}}
\newcommand{\Carver}{\operatorname{Carver}}
\newcommand{\Def}{\operatorname{Def}}
\newcommand{\sgn}{\operatorname{sgn}}
\DeclarePairedDelimiter{\abs}{\lvert}{\rvert}
\DeclarePairedDelimiter{\norm}{\lVert}{\rVert}

\title{Когда прямоугольник помещается в $n$-мерную коробку?\\ Явный критерий с формальным доказательством}
\author{Михаил Левит}
\date{29 сентября 2026~г.}

\begin{document}

\maketitle

\begin{abstract}
Мы даём необходимое и достаточное условие того, что прямоугольник $A\times B$ после произвольного
движения помещается в $n$-мерный прямоугольный параллелепипед (коробку) с рёбрами $p_1,\dots,p_n$, для
любого $n\ge2$. Условие --- конечная дизъюнкция: должны существовать две оси $k\ne l$ коробки и разбиение
остальных осей на множества $S$ и $T$ такие, что прямоугольник со сторонами
$\sqrt{A^2-\norm{p_S}^2}_+$ и $\sqrt{B^2-\norm{p_T}^2}_+$ должен помещаться в прямоугольник
$p_k\times p_l$, что проверяется плоским критерием Карвера. Геометрически каждая сторона прямоугольника
сначала идёт вдоль диагонали координатного подпространства, а остатки укладываются в одну координатную
плоскость. Тем самым решён случай $k=2$ общей задачи Джеррарда и Ветцеля (2004) о вложении
$k$-мерной коробки в $d$-мерную; при $n=3$ для единичного куба получается их Теорема~5. Доказательство
сводит задачу к ширинам прямоугольника вдоль осей и изучает минимизатор суммарного <<смешивания>> двух
сторон. Теорема проверена в системе доказательств Lean~4 с библиотекой Mathlib. Критерий конструктивен; на
Python он решает вопрос о вложении за $6$--$200$ микросекунд при $n=3,\dots,6$.
\end{abstract}

\section{Введение}

Когда одна коробка помещается в другую? Для отрезка длины $A$ и коробки с рёбрами $p_1,\dots,p_n$ ответ
очевиден: $A^2\le p_1^2+\dots+p_n^2$. Джеррард и Ветцель \cite[с.~22]{JerrardWetzel2004} называют общий
вопрос --- найти необходимые и достаточные условия того, что $k$-мерная коробка помещается в $d$-мерную ---
<<общей нерешённой задачей о вложении коробок в старших размерностях>> и отмечают, что при $k=1$ она
тривиальна. Первый нетривиальный случай --- прямоугольник, $k=2$. На плоскости ($d=2$) его решил Карвер
\cite{Carver1957} (см. также Ветцель \cite{Wetzel2000}). Для единичного куба ($d=3$) наибольший квадрат
имеет сторону $3\sqrt2/4$ (Ньивланд, опубликовано ван Свинденом \cite{vanSwinden1816}; задача принца
Руперта), а Джеррард и Ветцель \cite[Theorem~5]{JerrardWetzel2004} нашли наибольший прямоугольник
$L\times\lambda L$ в единичном кубе. Аналогичная задача для трёхмерной коробки в трёхмерной коробке решена в
\cite{Levit2026}. До сих пор точный ответ для прямоугольника был известен лишь для плоскости (Карвер) и для
единичного куба (Джеррард и Ветцель).

В этой работе такой критерий дан для любого $n\ge2$ (теорема~\ref{thm:main}). Это дизъюнкция
$\tfrac12 n(n-1)2^{n-2}$ явных проверок, каждая из которых --- плоский тест Карвера. Теорема говорит, что
оптимальное положение устроено просто: каждая сторона прямоугольника делится на часть, идущую вдоль
диагонали координатного подпространства (натянутого на оси $S$ для стороны $A$ и на оси $T$ для
стороны $B$), и остаток, лежащий в одной координатной плоскости $(k,l)$; там два остатка укладываются как
плоский прямоугольник в прямоугольник $p_k\times p_l$.

Доказательство (разделы~\ref{sec:suff}--\ref{sec:I}) элементарно. После обычного сведения к ширинам
прямоугольника вдоль осей (раздел~\ref{sec:reform}) задача становится системой для двух единичных векторов
$u$, $v$ с одним условием ортогональности. Допустимые конфигурации описываются двумя числами $(L,D)$,
лежащими под <<шатром>>; берётся конфигурация с наименьшим $D$ --- суммарным смешиванием сторон. Три
леммы показывают, что у такой конфигурации либо нет смешанных осей, либо их ровно две, и в обоих
случаях критерий выполнен.

Вся теорема проверена в системе доказательств Lean~4 с библиотекой Mathlib (раздел~\ref{sec:lean}).
Формальная формулировка использует только определение коробки в $\R^n$ и вложения изометрией;
доказательство не использует аксиом сверх стандартных. Критерий конструктивен: если прямоугольник
помещается, он даёт явное положение (раздел~\ref{sec:suff}), а реализация на Python на два--пять порядков
быстрее численного поиска по поворотам, который к тому же вблизи границы часто не находит существующего
положения (раздел~\ref{sec:comp}).

\section{Теорема}\label{sec:thm}

\subsection{Обозначения и критерий Карвера}

\emph{Коробка} --- это $P=[0,p_1]\times\dots\times[0,p_n]\subset\R^n$, $p_i\ge0$, $n\ge2$.
\emph{Прямоугольник} $A\times B$, $A,B\ge0$, --- множество $\{x_0+sAu+rBv:\ s,r\in[0,1]\}$ для некоторых
$x_0\in\R^n$ и ортонормированной пары $u,v\in\R^n$; он \emph{помещается} в $P$, если $x_0,u,v$ можно выбрать
так, чтобы он лежал в $P$. Отражения разрешены (это ничего не меняет: прямоугольник зеркально
симметричен). Для $S\subseteq[n]=\{1,\dots,n\}$ положим $\norm{p_S}^2=\sum_{i\in S}p_i^2$
($\norm{p_\varnothing}=0$) и $x_+=\max(x,0)$; пишем $\sqrt{x}_+$ вместо $\sqrt{x_+}$.

\medskip\noindent\textbf{Критерий Карвера} \cite{Carver1957} в форме Ветцеля \cite[Theorem~1]{Wetzel2000}.
Пусть $X\ge Y\ge0$ --- стороны плоского прямоугольника (контейнера), $a\ge b\ge0$ --- стороны другого.
Прямоугольник $a\times b$ помещается в $X\times Y$ тогда и только тогда, когда
\begin{enumerate}
\item[(i)] $a\le X$ и $b\le Y$, или
\item[(ii)] $a>X$, $b\le Y$ и $\displaystyle\Bigl(\frac{X+Y}{a+b}\Bigr)^2+\Bigl(\frac{X-Y}{a-b}\Bigr)^2\ge2$.
\end{enumerate}
Через $\Carver(a,b;X,Y)$ обозначаем это условие после сортировки обеих пар по убыванию. В случае~(ii)
$a>X\ge Y\ge b$, так что деления на ноль нет. При $b=0$ условие~(ii) превращается в $a^2\le X^2+Y^2$ ---
критерий для отрезка. Нужна только монотонность: если прямоугольник $a\times b$ помещается в $X\times Y$,
то помещается и любой $a_1\times b_1$ при $a_1\le a$, $b_1\le b$, в любой $X_1\times Y_1$ при $X_1\ge X$,
$Y_1\ge Y$.

\subsection{Формулировка}

\begin{theorem}\label{thm:main}
Пусть $n\ge2$, $p\in\R^n_{\ge0}$ и $A,B\ge0$. Прямоугольник $A\times B$ помещается в коробку $p$ тогда и
только тогда, когда
\begin{itemize}
\item[\textbf{(II)}] существуют оси $k\ne l$ и разбиение $S\sqcup T=[n]\setminus\{k,l\}$ такие, что
\[
\Carver(a',b';\,p_k,p_l),\qquad a'=\sqrt{A^2-\norm{p_S}^2}_+,\quad b'=\sqrt{B^2-\norm{p_T}^2}_+ .
\]
\end{itemize}
\end{theorem}

Порядок $a'\ge b'$ не гарантирован даже при $A\ge B$ (при $A=10$, $B=9$, $\norm{p_S}=9$, $T=\varnothing$
получается $a'=\sqrt{19}<b'=9$); поэтому $\Carver$ сортирует аргументы. Условие (II) не меняется при
перестановке $k$ и $l$, так что в нём $\tfrac12 n(n-1)2^{n-2}$ проверок: $1$, $6$, $24$, $80$, $240$ при
$n=2,\dots,6$. При $n=2$ это сам критерий Карвера.

Смысл (II) виден из построения в разделе~\ref{sec:suff}: сторона $A$ идёт, сколько может, вдоль
диагонали $p_S$ координатного подпространства осей $S$, сторона $B$ --- так же вдоль $p_T$, а остатки
$a'$, $b'$ укладываются по Карверу в плоскость осей $k,l$. Следующий частный случай (II) (см.
раздел~\ref{sec:I}) не требует плоскости вовсе:
\begin{itemize}
\item[\textbf{(I)}] существуют непересекающиеся $S,T\subseteq[n]$ с $A\le\norm{p_S}$ и $B\le\norm{p_T}$.
\end{itemize}

\begin{example}\label{ex:cube}
Квадрат $A=B$ в единичном кубе $p=(1,1,1)$. Возьмём $(k,l)=(1,2)$, $S=\{3\}$, $T=\varnothing$. Тогда
$a'=\sqrt{A^2-1}$, $b'=A$, и при $A>1$ условие Карвера (ii) в единичном квадрате имеет вид
$(2/(A+\sqrt{A^2-1}))^2\ge2$, то есть $A+\sqrt{A^2-1}\le\sqrt2$, или $A\le3\sqrt2/4\approx1{,}0607$.
Остальные проверки не дают большего, так что это наибольший квадрат в единичном кубе --- в согласии с
результатом Ньивланда. На рис.~\ref{fig:example} показано положение из раздела~\ref{sec:suff}: вершины
$(\tfrac14,0,0)$, $(0,\tfrac14,1)$, $(\tfrac34,1,1)$, $(1,\tfrac34,0)$; сторона $v$ лежит в плоскости осей
$1,2$, а у стороны $u$ компонента вдоль оси $3\in S$ равна $A\,u_3=1=p_3$.
\end{example}

\begin{figure}[ht]
\centering
\begin{tikzpicture}[scale=3.0,
  x={(1cm,0cm)}, y={(0.5cm,0.36cm)}, z={(0cm,1cm)}, line join=round]
  % куб
  \draw[gray] (0,0,0)--(1,0,0)--(1,1,0)--(0,1,0)--cycle;
  \draw[gray] (0,0,1)--(1,0,1)--(1,1,1)--(0,1,1)--cycle;
  \foreach \x/\y in {0/0,1/0,1/1,0/1} \draw[gray] (\x,\y,0)--(\x,\y,1);
  \draw[->] (1,0,0)--(1.3,0,0) node[right] {$1$};
  \draw[->] (1,1,0)--(1,1.4,0) node[right] {$2$};
  \draw[->] (0,0,1)--(0,0,1.3) node[above] {$3$};
  % квадрат
  \fill[blue!25,opacity=0.8] (0.25,0,0)--(0,0.25,1)--(0.75,1,1)--(1,0.75,0)--cycle;
  \draw[blue!70!black,thick] (0.25,0,0)--(0,0.25,1)--(0.75,1,1)--(1,0.75,0)--cycle;
  \foreach \P in {(0.25,0,0),(0,0.25,1),(0.75,1,1),(1,0.75,0)} \fill[blue!70!black] \P circle (0.5pt);
  \node[below] at (0.25,0,0) {\small $(\frac14,0,0)$};
  \node[left,xshift=-4pt,yshift=2pt] at (0,0.25,1) {\small $(0,\frac14,1)$};
  \node[above,yshift=1pt] at (0.75,1,1) {\small $(\frac34,1,1)$};
  \node[right,xshift=2pt] at (1,0.75,0) {\small $(1,\frac34,0)$};
\end{tikzpicture}
\hspace{1.2cm}
\begin{tikzpicture}[scale=3.0]
  \draw[gray] (0,0) rectangle (1,1);
  \draw[->] (1,0)--(1.25,0) node[right] {$1$};
  \draw[->] (0,1)--(0,1.25) node[above] {$2$};
  \fill[blue!25] (0.25,0)--(1,0.75)--(0.75,1)--(0,0.25)--cycle;
  \draw[blue!70!black,thick] (0.25,0)--(1,0.75)--(0.75,1)--(0,0.25)--cycle;
  \node[below right,xshift=-1pt] at (0.625,0.375) {\small $b'$};
  \draw[->,shorten >=1pt] (-0.12,-0.12) node[below left,inner sep=1pt] {\small $a'$}--(0.12,0.12);
  \node[below] at (0.5,0) {\small $p_1=1$};
  \node[left] at (0,0.5) {\small $p_2=1$};
\end{tikzpicture}
\caption{Слева: наибольший квадрат $A=B=3\sqrt2/4$ в единичном кубе в положении, построенном в
разделе~\ref{sec:suff} для $(k,l)=(1,2)$, $S=\{3\}$, $T=\varnothing$. Справа: его проекция на плоскость осей
$1,2$ --- прямоугольник $a'\times b'$, $a'=\sqrt{A^2-1}=\sqrt2/4$, уложенный под углом $45^\circ$ по критерию
Карвера; остальная часть стороны $A$ идёт вдоль оси $3$.}
\label{fig:example}
\end{figure}

\section{Достаточность}\label{sec:suff}

Пусть выполнено (II). Критерий Карвера даёт ортонормированные $f_a,f_b\in\operatorname{span}(e_k,e_l)$ с
$a'\abs{(f_a)_i}+b'\abs{(f_b)_i}\le p_i$ при $i\in\{k,l\}$. Положим $\lambda=\min(1,A/\norm{p_S})$ и $\mu=\min(1,B/\norm{p_T})$ (при $\norm{p_S}=0$ полагаем $\lambda=0$, при $\norm{p_T}=0$ --- $\mu=0$), $p_S=\sum_{i\in S}p_ie_i$ и
\[
A\,u=\lambda\,p_S+a'f_a,\qquad B\,v=\mu\,p_T+b'f_b .
\]
\begin{itemize}
\item Нормы: $\norm{Au}^2=\lambda^2\norm{p_S}^2+a'^2=A^2$ и при $A\ge\norm{p_S}$, и при $A<\norm{p_S}$
(тогда $a'=0$). Аналогично $\norm{Bv}=B$.
\item Ортогональность: носители $S$, $T$ и $\{k,l\}$ попарно не пересекаются, $f_a\perp f_b$, поэтому $u\perp v$.
\item Ширины: на $S$ имеем $A\abs{u_i}=\lambda p_i\le p_i$ и $v_i=0$; на $T$ --- симметрично; на $\{k,l\}$ ---
по Карверу.
\end{itemize}
По разделу~\ref{sec:widths} это положение прямоугольника в коробке. (При $A=0$ положим $u=f_a$, при $B=0$ ---
$v=f_b$.) \qed

\section{Переформулировка}\label{sec:reform}

\subsection{Сведение к ширинам}\label{sec:widths}
Коробка --- пересечение слоёв $\{0\le x_i\le p_i\}$, сдвиги по разным координатам независимы. Поэтому
выпуклое тело можно сдвинуть внутрь $P$ тогда и только тогда, когда его ширина вдоль каждого $e_i$ не больше
$p_i$. Ширина прямоугольника $(Au,Bv)$ вдоль $e_i$ равна $A\abs{u_i}+B\abs{v_i}$. Итак,
\begin{equation}\label{eq:widths}
\text{прямоугольник помещается}\iff\exists\ \text{ортонормированные } u,v:\quad A\abs{u_i}+B\abs{v_i}\le p_i\ \ \text{для всех } i .
\end{equation}

Случай $B=0$ (отрезок): помещается тогда и только тогда, когда $A\le\norm{p}$; необходимость следует из
$A^2=\sum A^2u_i^2\le\sum p_i^2$. Это (I) при $S=[n]$, $T=\varnothing$. Дальше $A,B>0$. Оси с $p_i=0$ сразу
выбрасываем: на них $u_i=v_i=0$. Если для оставшихся осей найдены $S,T,k,l$ из (I) или (II), каждую
выброшенную ось добавляем в $S$; это только увеличивает $\norm{p_S}$, значит, уменьшает $a'$ и сохраняет
условие Карвера, а в (I) сохраняет $A\le\norm{p_S}$.

Если осей с $p_i>0$ меньше двух, то при $A,B>0$ прямоугольник не помещается (два ортонормированных вектора
нельзя уложить в одну координату), и (II) тоже ложно. Действительно, пусть $m$ --- единственная ось с $p_m>0$
(или таких осей нет). Если $m\in\{k,l\}$, то $a'=A>0$ и $b'=B>0$, а коробка $(p_m,0)$ вмещает только отрезок;
если $m\in S$, то $b'=B>0$, а коробка $(0,0)$ --- точка; если $m\in T$ (или $m$ нет), то $a'=A>0$, и коробка
снова точка. Так что теорема верна и в этом случае.

\subsection{Ослабление}\label{sec:relax}
В \eqref{eq:widths} можно требовать лишь $\norm u\ge1$, $\norm v\ge1$, $\langle u,v\rangle=0$: сжатие $u$
или $v$ в $\le1$ раз сохраняет ортогональность и неравенства. Обозначим $P_i=p_i^2$,
$\alpha_i=A\abs{u_i}/p_i$, $\beta_i=B\abs{v_i}/p_i$ ($\alpha_i+\beta_i\le1$) и
$\sigma_i=\sgn(u_iv_i)\in\{\pm1\}$ (любой знак при $u_iv_i=0$). Условия принимают вид
\[
\sum P_i\alpha_i^2\ge A^2,\qquad \sum P_i\beta_i^2\ge B^2,\qquad \sum\sigma_iP_i\alpha_i\beta_i=0 .
\]

\subsection{Активные оси}\label{sec:active}
Пусть $\rho_i=\alpha_i+\beta_i\in[0,1]$ и $p'_i=\rho_ip_i=A\abs{u_i}+B\abs{v_i}\le p_i$. В коробке $p'$ при
тех же $u,v$ каждая ось \emph{активна}: $\alpha'_i+\beta'_i=1$. Положим $\alpha'_i=(1+t_i)/2$,
$\beta'_i=(1-t_i)/2$, $t_i\in[-1,1]$. Условия раздела~\ref{sec:relax} для $p'$ те же, так как
$P'_i\alpha'^2_i=A^2u_i^2$ и т.\,д. Оси с $\rho_i=0$ выбрасываем.

Условия (I) и (II) монотонны по $p$: при росте $p$ нормы $\norm{p_S}$ растут, значит, $a'$, $b'$ не растут, а
коробка $(p_k,p_l)$ растёт. Поэтому (I)$\vee$(II) для $p'$ влечёт (I)$\vee$(II) для $p$.
\emph{Итак, достаточно доказать, что у активной конфигурации в коробке $p$ выполнено (I) или (II).}
Дальше, не меняя обозначений, через $p$ обозначаем коробку, в которой все оси активны.

Активная конфигурация --- пара $(t,\sigma)$, $t\in[-1,1]^n$, $\sigma\in\{\pm1\}^n$, для которой
$A\abs{u_i}=p_i(1+t_i)/2$, $B\abs{v_i}=p_i(1-t_i)/2$ и $\sgn(u_iv_i)=\sigma_i$. \emph{Стороны}:
$+=\{i:\sigma_i=+1\}$ и $-=\{i:\sigma_i=-1\}$. Ось \emph{смешанная}, если $\abs{t_i}<1$, то есть $u_iv_i\ne0$.

Обратно, пусть все $p_i>0$ и $(t,\sigma)$ --- любая пара. Положим $\abs{u_i}=p_i(1+t_i)/(2A)$,
$\abs{v_i}=p_i(1-t_i)/(2B)$ и выберем знаки так, чтобы $\sgn(u_iv_i)=\sigma_i$ при $u_iv_i\ne0$. Тогда все
ширины равны $p_i$ и
\[
\norm u^2=\sum\frac{P_i(1+t_i)^2}{4A^2},\qquad \norm v^2=\sum\frac{P_i(1-t_i)^2}{4B^2},\qquad
\langle u,v\rangle=\sum\frac{\sigma_iP_i(1-t_i^2)}{4AB}.
\]
Назовём $(t,\sigma)$ \emph{допустимой}, если $\norm u\ge1$, $\norm v\ge1$ и $\langle u,v\rangle=0$. По
разделу~\ref{sec:relax} допустимая конфигурация даёт положение прямоугольника. Все конфигурации, которые
строятся в разделе~\ref{sec:lemmas}, задаются парами $(t',\sigma')$, и их допустимость проверяется только
через эти три условия; их равносильная запись дана ниже.

\subsection{Шатёр}\label{sec:tent}
Пусть $L=\sum P_it_i$, $\Def_\pm=\sum_{i\in\pm}P_i(1-t_i^2)$, $W=\sum P_i$, $c_A=2A^2-W$, $c_B=W-2B^2$.
\begin{itemize}
\item Ортогональность $\sum\sigma_iP_i(1-t_i^2)=0$ означает $\Def_+=\Def_-=:D$.
\item Тогда $\sum P_it_i^2=W-2D$, и
$\sum P_i(1+t_i)^2\ge4A^2\iff L-D\ge c_A$, $\ \sum P_i(1-t_i)^2\ge4B^2\iff L+D\le c_B$.
\end{itemize}
Итак, конфигурация с $\Def_+=\Def_-=D$ допустима тогда и только тогда, когда точка $(L,D)$ лежит под
<<шатром>> $\{D\le L-c_A,\ D\le c_B-L\}$ (рис.~\ref{fig:tent}). Шатёр 1-липшицев, поэтому:

\medskip\noindent$(\star)$ \emph{Если $(t,\sigma)$ допустима, а $(t',\sigma')$ --- другая конфигурация с
$\Def'_+=\Def'_-=D'$ и $D'+\abs{L'-L}\le D$, то и $(t',\sigma')$ допустима.}
Действительно, $L'-D'\ge L'-D+\abs{L'-L}\ge L-D\ge c_A$ и $L'+D'\le L+D\le c_B$.

\medskip
Смешанная ось даёт положительный вклад в $\Def$ своей стороны. Так как $\Def_+=\Def_-$, смешанные оси либо
есть на обеих сторонах, либо их нет вовсе; ровно одна смешанная ось невозможна.

\begin{figure}[ht]
\centering
\begin{tikzpicture}[scale=1.0]
  \fill[blue!12] (-1,0)--(1,2)--(3,0)--cycle;
  \draw[->] (-2.2,0)--(4.4,0) node[right] {$L$};
  \draw[->] (-2,0)--(-2,2.9) node[above] {$D$};
  \draw[thick] (-1.2,-0.2)--(1.6,2.6) node[pos=0.74,left,xshift=-3pt] {\small $D=L-c_A$};
  \draw[thick] (0.4,2.6)--(3.2,-0.2) node[pos=0.26,right,xshift=3pt] {\small $D=c_B-L$};
  \node[below,xshift=-14pt] at (-1,0) {\small $c_A$};
  \node[below,xshift=14pt] at (3,0) {\small $c_B$};
  \fill[blue!30] (1.2,0.8)--(0.4,0)--(2.0,0)--cycle;
  \draw[dashed] (0.4,0)--(1.2,0.8)--(2.0,0);
  \fill (1.2,0.8) circle (1.6pt) node[above] {\small $(L,D)$};
\end{tikzpicture}
\caption{Допустимым конфигурациям отвечают точки $(L,D)$ под шатром. Каждая точка $(L',D')$ с
$D'+\abs{L'-L}\le D$ (тёмный треугольник под $(L,D)$) тоже лежит под шатром; это $(\star)$.}
\label{fig:tent}
\end{figure}

\section{Три леммы}\label{sec:lemmas}

\begin{lemma}[C: не больше двух смешанных осей]\label{lem:C}
Если у допустимой конфигурации нет смешанных осей, то выполнено (I). Если смешанных осей ровно две,
$k\in+$ и $l\in-$, то выполнено (II) с этими $k,l$.
\end{lemma}
\begin{proof}
Нет смешанных осей. Положим $S=\{t=+1\}$, $T=\{t=-1\}$. Тогда $D=0$, $L=\norm{p_S}^2-\norm{p_T}^2$,
$W=\norm{p_S}^2+\norm{p_T}^2$, и шатёр даёт $\norm{p_S}^2\ge A^2$, $\norm{p_T}^2\ge B^2$; это (I).

Две смешанные оси $k\in+$, $l\in-$. На остальных осях $u_iv_i=0$, поэтому для $U=(u_k,u_l)$, $V=(v_k,v_l)$
выполнено $\langle U,V\rangle=\langle u,v\rangle=0$; кроме того, $U,V\ne0$, так как на смешанной оси
$u_iv_i\ne0$. Значит, прямоугольник $A\abs U\times B\abs V$ с направлениями $U/\abs U$, $V/\abs V$ лежит в
2-коробке $(p_k,p_l)$: его ширины $A\abs{u_i}+B\abs{v_i}=p_i$. Пусть $S=\{t=+1\}$; там $v_i=0$ и
$A\abs{u_i}=p_i$. Тогда $A^2\abs U^2=A^2\norm u^2-\norm{p_S}^2\ge A^2-\norm{p_S}^2$ (оси с $t=-1$ не дают
вклада в $u$), то есть $A\abs U\ge a'$. Аналогично $B\abs V\ge b'$ при $T=\{t=-1\}$. По монотонности
$\Carver(a',b';p_k,p_l)$; это (II).
\end{proof}

\begin{lemma}[R: разные знаки]\label{lem:R}
Пусть у всех смешанных осей стороны $+$ выполнено $t\ge0$, а у всех смешанных осей стороны $-$ выполнено
$t\le0$ (или наоборот). Тогда выполнено (I).
\end{lemma}
\begin{proof}
Округлим: $t\mapsto+1$ у смешанных осей стороны $+$ и $t\mapsto-1$ у смешанных осей стороны $-$.
\begin{itemize}
\item Сдвиг $L$ от стороны $+$: $\Delta_+=\sum P(1-t)\in[0,D]$, так как $1-t\le1-t^2$ при $t\in[0,1]$.
\item Сдвиг $L$ от стороны $-$: $\Delta_-=-\sum P(1+t)\in[-D,0]$, так как $1+t\le1-t^2$ при $t\in[-1,0]$.
\item У новой конфигурации $\Def'_+=\Def'_-=0$ и $\abs{L'-L}=\abs{\Delta_++\Delta_-}\le D$, поскольку
слагаемые разных знаков.
\end{itemize}
По $(\star)$ она допустима. Смешанных осей у неё нет, так что по лемме~\ref{lem:C} выполнено (I). Случай
<<наоборот>> симметричен.
\end{proof}

\begin{lemma}[V: ход по кривой]\label{lem:V}
Пусть $i\ne j$ --- смешанные оси одной стороны (скажем, $+$), $k$ --- смешанная ось другой стороны,
$s_0=t_k$. Пусть не выполнено <<$t_i=t_j$ и $t_is_0\le0$>>, то есть либо (a) $t_i\ne t_j$, либо (b) $t_i=t_j$ и
$t_is_0>0$. Тогда существует допустимая конфигурация с меньшим $D$.
\end{lemma}
\begin{proof}
Меняем только $(t_i,t_j,t_k)$, сохраняя $L$ и $\Def_+-\Def_-$:
\[
P_it_i+P_jt_j=\lambda(s):=L_0-P_ks,\qquad P_it_i^2+P_jt_j^2=\mu(s):=K_0+P_ks^2,\qquad t_k=s .
\]
Тогда $D=\Def_-=\mathrm{const}+P_k(1-s^2)$ строго убывает с ростом $\abs s$, а $L$ не меняется. По
$(\star)$ конфигурация остаётся допустимой, пока $(t_i,t_j,t_k)\in[-1,1]^3$. (Остальные смешанные оси стороны
$+$, если они есть, входят в константы $L_0$, $K_0$.)

Прямая $\{P_iX+P_jY=\lambda\}$ и эллипс $\{P_iX^2+P_jY^2=\mu\}$ пересекаются тогда и только тогда, когда
$h(s):=\mu(s)-\lambda(s)^2/(P_i+P_j)\ge0$: по неравенству Коши--Буняковского минимум $P_iX^2+P_jY^2$ на прямой
достигается при $X=Y=\tau:=\lambda/(P_i+P_j)$, а $h=0$ означает касание, $t_i=t_j=\tau$. При этом
$h'(s)=2P_k(s+\tau)$.

(a) Если $t_i\ne t_j$, якобиан отображения $(X,Y)\mapsto(P_iX+P_jY,\,P_iX^2+P_jY^2)$ равен
$2P_iP_j(Y-X)\ne0$, и по теореме о неявной функции решение продолжается на $s$, близкие к $s_0$, в частности на
$\abs s>\abs{s_0}$ (при $s_0=0$ --- в любую сторону). Точка $(t_i,t_j)$ лежит внутри квадрата $(-1,1)^2$ и
$\abs{s_0}<1$, так что при малом ходе всё остаётся в $[-1,1]$, и $D$ строго убывает.

(b) Если $t_i=t_j=\tau$ и $\tau s_0>0$, то $s_0\ne0$, $h(s_0)=0$, а $h'(s_0)=2P_k(s_0+\tau)$ имеет знак $s_0$.
Значит, $h>0$ при малом росте $\abs s$, и точек пересечения две. Так как $i,j$ смешанные, $\abs\tau<1$; выберем
$\varepsilon>0$ с $\abs\tau+2\varepsilon<1$. Обе точки пересечения непрерывно зависят от $s$ и при $s\to s_0$
стремятся к $(\tau,\tau)$ (расстояние до точки касания порядка $\sqrt h$). Поэтому при достаточно малом
изменении $s$ обе координаты остаются в $(-\abs\tau-\varepsilon,\abs\tau+\varepsilon)\subset(-1,1)$, и
$\abs s<1$. Снова $D$ строго убывает.
\end{proof}

\section{Необходимость}\label{sec:nec}

По разделу~\ref{sec:reform} достаточно рассмотреть допустимую активную конфигурацию в коробке $p$ со всеми
$p_i>0$. При каждом из конечного числа $\sigma$ множество допустимых $t\in[-1,1]^n$ задаётся замкнутыми условиями
$\Def_+=\Def_-$, $L-D\ge c_A$, $L+D\le c_B$, где $L$, $\Def_\pm$ и $D$ --- многочлены от $t$; поэтому оно
компактно. При $\sigma$ исходной конфигурации оно непусто, а $D\ge0$ непрерывна, значит, существует конфигурация с \emph{наименьшим}~$D$.
\begin{enumerate}
\item Если на какой-то стороне есть две смешанные оси с разными $t$, то на другой стороне тоже есть смешанная
ось (раздел~\ref{sec:tent}), и лемма~\ref{lem:V}(a) уменьшает $D$ --- противоречие. Значит, у всех смешанных осей
стороны $+$ общее $t=\tau$, а у всех смешанных осей стороны $-$ общее $t=s$.
\item Если на каждой стороне не больше одной смешанной оси, то по разделу~\ref{sec:tent} их $0$ или $1+1$, и
работает лемма~\ref{lem:C}.
\item Иначе на какой-то стороне (скажем, $+$) не меньше двух смешанных осей с общим $\tau$. Возьмём смешанную ось
$k$ другой стороны, $t_k=s$. Если $\tau s>0$, лемма~\ref{lem:V}(b) уменьшает $D$ --- противоречие. Значит,
$\tau s\le0$, выполнены условия знаков леммы~\ref{lem:R}, и получаем (I). \qed
\end{enumerate}

\begin{corollary}\label{cor:structure}
Пусть (I) неверно. Тогда у любой допустимой конфигурации с наименьшим $D$ ровно две смешанные оси, по одной на
каждой стороне, и их $t$ одного знака. Иными словами, оптимальное положение имеет вид (II) с поворотом в одной
координатной плоскости.
\end{corollary}
\begin{proof}
Отсутствие смешанных осей дало бы (I) по лемме~\ref{lem:C}; две и более на одной стороне по пп.~1 и 3 ведут к
противоречию или к (I); разные знаки дали бы (I) по лемме~\ref{lem:R}.
\end{proof}

\section{(I) --- частный случай (II)}\label{sec:I}

Пусть $A,B>0$ и выполнено (I). Тогда $S$ и $T$ непусты, так как $\norm{p_S}\ge A>0$ и $\norm{p_T}\ge B>0$.
Пусть $R=[n]\setminus(S\cup T)$; возьмём $k\in S$, $l\in T$ и положим $S'=(S\setminus\{k\})\cup R$,
$T'=T\setminus\{l\}$, так что $S'\sqcup T'=[n]\setminus\{k,l\}$. Тогда
$\norm{p_{S'}}^2=\norm{p_S}^2-p_k^2+\norm{p_R}^2\ge A^2-p_k^2$, откуда $a'^2\le p_k^2$; аналогично
$b'^2\le p_l^2$, и условие Карвера выполнено без поворота. При $B=0$ возьмём любые $k\ne l$, $S$ --- остальные
оси, $T=\varnothing$: отрезок $a'$ помещается в $(p_k,p_l)$ тогда и только тогда, когда
$a'\le\sqrt{p_k^2+p_l^2}$, то есть $A\le\norm p$.

Вместе с разделами~\ref{sec:suff}--\ref{sec:nec} это доказывает теорему~\ref{thm:main}. Попутно: у
\eqref{eq:widths} есть решение, в котором не больше двух координат с $u_iv_i\ne0$; и если (I) неверно, то у двух
смешанных осей в (II) $t$ одного знака.

\section{Единичный куб: Теорема~5 Джеррарда и Ветцеля}\label{sec:jw}

Джеррард и Ветцель \cite[Theorem~5, формула~(8)]{JerrardWetzel2004} нашли наибольший прямоугольник
$L\times\lambda L$, $0\le\lambda\le1$, в единичном кубе: $L_{\max}=f_1(\lambda)$ на $[0,\lambda_1]$,
$\sqrt2$ на $[\lambda_1,\lambda_2]$ и $f_3(\lambda)$ на $[\lambda_2,1]$, где
\[
f_1(\lambda)=\frac{3}{\sqrt{3-\lambda^2}+\sqrt2\lambda},\qquad f_3(\lambda)=\frac{3}{\sqrt2+\sqrt{3\lambda^2-1}},\qquad
\lambda_1=1-\tfrac{\sqrt2}{2},\quad \lambda_2=\tfrac{\sqrt2}{2}.
\]
Теорема~\ref{thm:main} даёт тот же ответ другим путём. Пусть $p=(1,1,1)$, $A=L$, $B=\lambda L$, и пусть $m$ ---
третья ось, $\{k,l,m\}=\{1,2,3\}$.
\begin{itemize}
\item $m\in S$, $A>1$: $a'=\sqrt{A^2-1}$, $b'=\lambda A$, и Карвер в единичном квадрате даёт
$\sqrt{A^2-1}+\lambda A\le\sqrt2$, то есть $(1-\lambda^2)A^2+2\sqrt2\lambda A-3\le0$, или $A\le f_1(\lambda)$.
\item $m\in T$, $B\le1$: $b'=0$, отрезок $A$ в единичном квадрате, $A\le\sqrt2$; здесь $B=\lambda\sqrt2\le1$
означает $\lambda\le\sqrt2/2$.
\item $m\in T$, $B>1$: $a'=A$, $b'=\sqrt{\lambda^2A^2-1}$, и Карвер даёт $A+\sqrt{\lambda^2A^2-1}\le\sqrt2$, то есть
$A\le f_3(\lambda)$. Это осуществимо только при $\lambda\ge\lambda_2$: левая часть растёт по $A$ и при $A=1/\lambda$
равна $1/\lambda$, так что нужно $1/\lambda\le\sqrt2$.
\item Прямоугольник без поворота даёт $\min(\sqrt2,1/\lambda)$, что нигде не лучше.
\end{itemize}
Знаменатели $f_1$ и $f_3$ растут по $\lambda$, так что $f_1$ и $f_3$ убывают, и $f_1(\lambda_1)=f_3(\lambda_2)=\sqrt2$.
На $[\lambda_2,1]$ выполнено $f_3\ge f_1$: разность знаменателей
$\varphi=(\sqrt{3-\lambda^2}+\sqrt2\lambda)-(\sqrt2+\sqrt{3\lambda^2-1})$ равна нулю при $\lambda=1$ и убывает, так как
$\varphi'\le\sqrt2-3\lambda/\sqrt{3\lambda^2-1}\le\sqrt2-3/\sqrt2<0$. Значит, максимум по этим вариантам --- в
точности $L_{\max}(\lambda)$. У Джеррарда и Ветцеля разбор идёт по положению вершин, здесь --- перебор
разбиений осей. Численно перебор всех $(k,l,S,T)$ с бисекцией совпадает с формулой~(8) на $1001$ значении
$\lambda\in[0,1]$ с точностью $3{,}3\cdot10^{-16}$.

\section{Формальное доказательство}\label{sec:lean}

Теорема доказана в Lean~4 с Mathlib; исходные тексты --- в каталоге \texttt{lean/} репозитория \cite{Repo}.
Смысл главной теоремы зависит только от следующих определений (файл \texttt{Widths.lean}):
\begin{verbatim}
abbrev En (n : ℕ) := EuclideanSpace ℝ (Fin n)
def boxN (q : Fin n → ℝ) : Set (En n) := {x | ∀ i, 0 ≤ x i ∧ x i ≤ q i}
noncomputable def pad (h : m ≤ n) (q : Fin m → ℝ) : Fin n → ℝ :=
  Function.extend (Fin.castLE h) q 0
def FitsN (p q : Fin n → ℝ) : Prop := ∃ f : En n ≃ᵢ En n, f >> boxN q ⊆ boxN p
\end{verbatim}
Здесь \texttt{pad h ![A, B]} --- вектор рёбер $(A,B,0,\dots,0)$, так что прямоугольник --- это коробка
\texttt{boxN (pad h ![A, B])} в $\R^n$, а \texttt{FitsN} требует изометрии $\R^n$ (любого движения, включая
отражения), переводящей одну коробку в другую. Критерий Карвера записан дословно как в разделе~\ref{sec:thm}
(\texttt{CarverS} для отсортированных сторон, \texttt{Carver} сортирует через \texttt{max}/\texttt{min}).
Главная теорема (файл \texttt{RectCarver.lean}):
\begin{verbatim}
theorem fitsN_rect_iff_explicit (h : 2 ≤ n) (p : Fin n → ℝ) (hp : ∀ i, 0 ≤ p i)
    {A B : ℝ} (hA : 0 ≤ A) (hB : 0 ≤ B) :
    FitsN p (pad h ![A, B]) ↔
      ∃ k l : Fin n, k ≠ l ∧ ∃ S T : Finset (Fin n), Disjoint S T ∧
        S ∪ T = univ \ {k, l} ∧
        Carver (√(A ^ 2 - ∑ i ∈ S, p i ^ 2)) (√(B ^ 2 - ∑ i ∈ T, p i ^ 2))
          (p k) (p l)
\end{verbatim}
(в Mathlib $\sqrt x=0$ при $x\le0$, так что срезка $(\cdot)_+$ получается сама). Команда
\texttt{\#print axioms} показывает только стандартные аксиомы \texttt{propext}, \texttt{Classical.choice},
\texttt{Quot.sound}; \texttt{sorry} в исходниках нет. Версии: Lean~4 v4.34.0-rc2, Mathlib v4.34.0-rc2.
Гипотезы --- только $n\ge2$, $p_i\ge0$, $A,B\ge0$, так что вырожденные случаи $A=0$ или $B=0$ (отрезок,
точка) покрыты той же формулировкой.

Формализация следует разделам~\ref{sec:suff}--\ref{sec:I} с несколькими упрощениями. Вместо перехода к
активным осям она работает с $x=Au$, $y=Bv$ и локальными ширинами $r_i=\abs{x_i}+\abs{y_i}\le p_i$, которые
сохраняются всеми шагами; минимизирует $\sum\abs{x_iy_i}$ (это $D/2$) сразу по всем решениям ослабленной задачи,
не фиксируя $\sigma$; лемма~\ref{lem:V} доказана без теоремы о неявной функции, через явную параметризацию
прямой $X=\tau+P_j\theta$, $Y=\tau-P_i\theta$, для которой $P_iX^2+P_jY^2=(P_i+P_j)\tau^2+P_iP_j(P_i+P_j)\theta^2$.
Сам критерий Карвера выведен из леммы о полосе из \cite{Levit2026}: при $a>X$ оптимальное положение
прямоугольника $a\times b$ касается полосы ширины $X$, и остаётся сравнить его протяжённость вдоль полосы с
$Y$, что является чистой алгеброй.

\section{Вычисления}\label{sec:comp}

Критерий перебирает $\tfrac12n(n-1)2^{n-2}$ случаев, каждый из которых --- тест Карвера с двумя квадратными
корнями; число случаев растёт по $n$ экспоненциально, но при $n\le10$ их не больше $11\,520$. Далее, раздел~\ref{sec:suff} превращает успешный случай в явную ортонормированную пару $(u,v)$. Мы сравнили
реализацию на Python (\texttt{tools/rect\_fit.py} в \cite{Repo}) с двумя численными поисками положения:
методом Нелдера--Мида для $F(u,v)=\max_i(A\abs{u_i}+B\abs{v_i})/p_i$ по ортонормированным парам (как в
\cite{Levit2026}) и методом SLSQP для гладкой задачи <<минимизировать $t$ при $Aa_i+Bb_i\le tp_i$,
$a_i\ge\abs{u_i}$, $b_i\ge\abs{v_i}$, $\norm u=\norm v=1$, $u\perp v$>>, оба из случайных стартов. Численный поиск
отвечает <<помещается>>, как только какой-то старт достигает $F\le1$; подтвердить, что прямоугольник не
помещается, он не может. Для каждого $n$ взято $50$ пар у границы: прямоугольник масштабируется на
$t^*(1\pm e)$, где $t^*$ --- наибольший масштаб, при котором он помещается, и $10^{-4}\le e\le10^{-2}$;
помещается примерно половина. Время на пару (Python~3.10, одно ядро). Критерий останавливается на первом
удачном случае; столбец <<критерий>> усреднён по всем парам, включая не помещающиеся, для которых нужны все
случаи, а столбец <<критерий + положение>> --- только по помещающимся, поэтому он может быть меньше.

\begin{center}\small
\begin{tabular}{ccccc}
\toprule
$n$ & критерий & критерий + положение & Нелдер--Мид, 20 стартов & SLSQP, 20 стартов\\
\midrule
3 & 6\,мкс & 7\,мкс & 1,0 с; пропущено 3 из 25 & 26 мс; пропущено 1 из 25\\
4 & 26\,мкс & 21\,мкс & 1,9 с; пропущено 19 из 22 & 40 мс; пропущено 1 из 22\\
5 & 80\,мкс & 36\,мкс & 2,4 с; пропущено 24 из 24 & 46 мс; пропущено 2 из 24\\
6 & 194\,мкс & 42\,мкс & 3,3 с; пропущено 24 из 24 & 63 мс; пропущено 3 из 24\\
\bottomrule
\end{tabular}
\end{center}

\noindent <<Пропущено>> значит, что поиск ответил <<не помещается>> для прямоугольника, который помещается.
Ни разу поиск не нашёл положения там, где критерий говорит <<не помещается>>. Все положения, построенные по критерию, прошли проверку (ортонормированность и ширины) --- и здесь, и в
отдельной самопроверке реализации на $20\,000$ случайных пар при $n=2,\dots,7$ (без замера времени). Сравнения чисел с плавающей точкой точны лишь в точной арифметике: у границы округление может
изменить ответ.

\section{Замечания}\label{sec:remarks}

Рассуждение использует одну пару ортогональных столбцов $(u,v)$, то есть одно условие ортогональности,
которое принимает вид $\Def_+=\Def_-$. Поэтому задача укладывается в двумерную картинку $(L,D)$ с шатром. У
коробки с тремя ненулевыми рёбрами таких условий три; каждая пара столбцов даёт своё разбиение осей на
стороны, и ход леммы~\ref{lem:V}, сохраняющий одно условие, ломает два других. Это пояснение, а не
доказательство. Для трёхмерной коробки в трёхмерной коробке явный критерий известен \cite{Levit2026}; при
$k\ge3$ и $d\ge4$ задача Джеррарда и Ветцеля, насколько нам известно, остаётся открытой.

\begin{thebibliography}{99}
\bibitem{Carver1957} W.~B.~Carver, Solution of Problem E1225, \emph{Amer. Math. Monthly} 64 (1957) 114--116.
\bibitem{JerrardWetzel2004} R.~P.~Jerrard, J.~E.~Wetzel, Prince Rupert's rectangles, \emph{Amer. Math. Monthly} 111 (2004) 22--31.
\bibitem{Levit2026} М.~Левит, Когда коробка помещается в коробку? Явный критерий с формальным
доказательством, 2026, \url{https://github.com/MikhailLevit1987/box-in-box},
\href{https://doi.org/10.5281/zenodo.22974798}{doi:10.5281/zenodo.22974798}.
\bibitem{Repo} М.~Левит, Прямоугольник в $n$-мерной коробке: статья, формализация в Lean и программы, 2026,
\url{https://github.com/MikhailLevit1987/rectangle-in-n-box}.
\bibitem{vanSwinden1816} J.~H.~van Swinden, \emph{Grondbeginsels der Meetkunde}, 2nd ed., P.~den Hengst en Zoon, Amsterdam, 1816, с.~512--513, 608--610.
\bibitem{Wetzel2000} J.~E.~Wetzel, Rectangles in rectangles, \emph{Math. Mag.} 73 (2000) 204--211.
\end{thebibliography}

\end{document}
